\documentclass[11pt]{article}
\usepackage[a4paper,margin=18mm]{geometry}
\usepackage{fontspec}
\setmainfont{Latin Modern Roman}
\usepackage{amsmath,amssymb,amsthm,amscd,graphicx,color}
\usepackage{xurl}
\usepackage[unicode,hidelinks]{hyperref}
\allowdisplaybreaks
\setlength\emergencystretch{3em}
\DeclareMathOperator{\id}{id}
\DeclareMathOperator{\lcm}{lcm}
\DeclareMathOperator{\End}{End}
\DeclareMathOperator{\Ind}{Ind}
\DeclareMathOperator{\Wor}{Wor}
\DeclareMathOperator{\GrAut}{GrAut}
\DeclareMathOperator{\Id}{Id}
\DeclareMathOperator{\Der}{Der}
\DeclareMathOperator{\Hom}{Hom}
\DeclareMathOperator{\Aut}{Aut}
\DeclareMathOperator{\im}{im}
\DeclareMathOperator{\gr}{\mathtt{gr}}
\DeclareMathOperator{\GVect}{\Gamma\hyp\Vect}
\DeclareMathOperator{\hyp}{-}
\newcommand{\g}[1]{t_{#1}^{(-1)}}

\newtheorem{Theorem}{Theorem}[section]
\newtheorem{claim}[Theorem]{Claim}
\newtheorem{maintheorem}[Theorem]{Main Theorem}
\newtheorem{Corollary}[Theorem]{Corollary}
\newtheorem{Lemma}[Theorem]{Lemma}
\newtheorem{Proposition}[Theorem]{Proposition}
\newtheorem{problem}[Theorem]{Problem}
{\theoremstyle{definition}
\newtheorem{Definition}[Theorem]{Definition}
\newtheorem{Remark}[Theorem]{Remark}}



\begin{document}
\begin{center}\Large Group rings over Z[(1+\allowbreak{}i)/2] of G(m,p,n) and its mystic partner are isomorphic for n>=\allowbreak{}1, even m>=\allowbreak{}2 and p dividing m\end{center}
\noindent\textbf{Stable ID:} 2020\par
\noindent\textbf{Authors:} Yuri Bazlov; Arkady Berenstein\par
\noindent\textbf{Original work:} Mystic Reflection Groups\par
\noindent\textbf{Version:} Official SIGMA 10 (2014) article 040, DOI 10.3842/SIGMA.2014.040; journal PDF and native TeX acquired 2026-10-01, exact bytes pinned by SHA256\par
\noindent\textbf{Selected task scope:} Full MainTheorem2.8 with MainTheorem2.5 hypotheses n>=\allowbreak{}1,m>=\allowbreak{}2 even,p a positive divisor ofm: R G(m,p,n) is isomorphic to R mu(G) for R=\allowbreak{}Z[(1+\allowbreak{}i)/2], hence the complex group rings are isomorphic. This does not assert the groups themselves are isomorphic or an integral-Z group-ring isomorphism.\par
\noindent\textbf{Difficulty:} H2: The proof constructs an explicit automorphism of the monomial-group algebra from parity-dependent root-system cocycles and finite diagonal idempotent factors. Restricting its parameter to i preserves the specified integral coefficient ring and sends reflection generators to combinations of mystic generators, proving a group-ring isomorphism even when the underlying groups are not isomorphic.\par
\noindent\textbf{Editorial boundary note (not author text):} The complete original Main Theorem 2.8 over R=\allowbreak{}Z[(1+\allowbreak{}i)/2] is retained with the explicit J\_c automorphism, cocycle/product identities, action comparison, mystic-partner description and full faithful-action/Dedekind Appendix argument required by the construction. The selected assumptions are n>=\allowbreak{}1, even m>=\allowbreak{}2 and p dividing m. J\_i fixes the common diagonal subgroup and sends each symmetric transposition to the exact displayed linear combination of sigma\_i and its inverse; inverse J\_(i inverse) is explicitly constructed. The cocycle identities are the author bounded routine root-system checks after their displayed formula. No group-isomorphism classification, invariant-ring theorem or odd-m extension is selected.\par
\noindent\textbf{Source:} \url{https://sigma-journal.com/2014/040/}\quad DOI: \url{https://doi.org/10.3842/SIGMA.2014.040}\par
\noindent\textbf{License and attribution:} Copyright retained by the named authors. Published by SIGMA under CC BY 4.0: \url{https://creativecommons.org/licenses/by/4.0/}. Source: \url{https://sigma-journal.com/about.html}. Adaptation consists of standalone excerpt formatting, cover and numbering restoration; original author language and mathematical text are retained.\par
\medskip\hrule\medskip\noindent\textbf{Author text begins below.}\par
\setcounter{section}{1}
\setcounter{subsection}{0}
\setcounter{subsubsection}{0}
\setcounter{Theorem}{3}
\setcounter{equation}{0}
% Original source lines 141--322
\section{Main results}
\label{sect:main}

We start with a~new notion of ``mystical equivalence'' of group actions, which is crucial for what follows.

\begin{Definition}
Let $ \mathop{\triangleright}\nolimits\colon G \times \mathcal V\to \mathcal V$,
$\mathop{\triangleright}\nolimits\nolimits'\colon G' \times \mathcal V\to \mathcal V $ be faithful actions of f\/inite
groups~$G$, respectively $G'$, on a~complex vector space~$\mathcal V$.
We say that the actions $\mathop{\triangleright}\nolimits$ and $\mathop{\triangleright}\nolimits'$ are
\textit{mystically equivalent}, if
\begin{gather*}
\rho(e_G)=\rho'(e_{G'}),
\end{gather*}
where $\rho\colon \mathbb C G \to \End_\mathbb C \mathcal V$ and $\rho'\colon \mathbb C G'\to \End_\mathbb C\mathcal V$
are the algebra homomorphisms def\/ined by the actions, and $e_G$ denotes the element $\sum\limits_{g\in G}g$ of $\mathbb C G$.
\end{Definition}

Mystical equivalence of the actions of~$G$ and $G'$ is a~strengthening of the condition that the respective spaces
$\mathcal V^G$ and $\mathcal V^{G'}$ of invariants are equal, due to the following obvious result.
\begin{Lemma}%\label{lemma:projection}
If a~$G$-action and a~$G'$-action on $\mathcal V$ are mystically equivalent, then $\mathcal V^G=\mathcal V^{G'}$.
\end{Lemma}

We will use the lemma in the situation where $\mathcal V = S(V)$ where $V$ is a~vector space over $\mathbb C$ with
a~chosen basis $\{x_1,\ldots,x_n\}$.
Throughout the paper, $n\ge 2$.
Denote by $\mathbb G_n$ the group of monomial matrices on $V$, that is, matrices in $\mathrm{GL}_n(\mathbb C)$ with
exactly $n$ non-zero entries.
In other words,
\begin{gather*}
\mathbb G_n=(\mathbb C^\times)^n\rtimes \mathbb S_n.
\end{gather*}
Here $(\mathbb C^\times)^n$ is naturally identif\/ied with the group of diagonal matrices in $\mathrm{GL}_n(\mathbb C)$
and acts on~$V$ by scaling the basis $\{x_1,\ldots,x_n\}$.
The symmetric group $\mathbb S_n$ is identif\/ied with the group of permutation matrices and acts on~$V$ by permuting the
same basis.
Note that $\mathbb G_n$ is the normalizer of the torus $(\mathbb C^\times)^n$ in $\mathrm{GL}_n(\mathbb C)$.
In particular, $\mathbb S_n$ acts on $(\mathbb C^\times)^n$ by conjugation.
When we write $tw\in \mathbb G_n$, we will imply that $t\in (\mathbb C^\times)^n$ and $w\in \mathbb S_n$; every element
of $\mathbb G_n$ can be uniquely written in this way.

Clearly, $\mathbb G_n$ is generated by $s_1,\ldots,s_{n-1}$ and $t_j^{(\zeta)}$, $1\le j\le n$, $\zeta\in \mathbb
C^\times$, where
\begin{itemize}\itemsep=0pt
\item
$s_i\in \mathbb S_n$ is the permutation of $\{x_1,\ldots,x_n\}$ which swaps $x_i$ and $x_{i+1}$;
\item
$t_j^{(\zeta)}\in (\mathbb C^\times)^n$ maps $x_k$ to $\zeta^{\delta_{jk}} x_k$.
\end{itemize}
We f\/ind it very convenient to use the linear character
\begin{gather*}
\det \colon \ \mathbb G_n \to \mathbb C^\times
\end{gather*}
of $\mathbb G_n$, which is just the restriction of the determinant character of $\mathrm{GL}_n(\mathbb C)$ to $\mathbb
G_n$.
In particular,
\begin{gather*}
\det s_i = -1,
\qquad
\det w\in\{\pm1\}\text{ for }w\in \mathbb S_n,
\qquad
\det\left(t_1^{(\zeta_1)} t_2^{(\zeta_2)} \cdots t_n^{(\zeta_n)}\right)=\zeta_1\zeta_2\cdots \zeta_n.
\end{gather*}

{\sloppy Next, we introduce \textit{two} dif\/ferent faithful actions of $\mathbb G_n$ on $S(V)$ using the natural basis
$\big\{x_1^{k_1} \cdots x_n^{k_n} :  k_1,\ldots,k_n\in \mathbb Z_{\ge0}\big\}$ of~$S(V)$.
(At the moment, we are not using any multiplication on $S(V)$.)

}

\begin{Proposition}\label{prop:automorphism}\quad
\begin{enumerate}\itemsep=0pt
\item[$(a)$] There exist $($unique$)$ faithful actions $\mathop{\triangleright}\nolimits_+$, $\mathop{\triangleright}\nolimits_-$ of
$\mathbb G_n$ on $S(V)$ such that
\begin{gather*}
s_i \mathop{\triangleright}\nolimits_\pm x_1^{k_1}\cdots x_n^{k_n} = (\pm 1)^{k_ik_{i+1}} x_1^{k_1}\cdots
x_i^{k_{i+1}}x_{i+1}^{k_i}\cdots x_n^{k_n},
\\
t_j^{(\zeta)}\mathop{\triangleright}\nolimits_\pm x_1^{k_1}\cdots x_n^{k_n} = \zeta^{k_j} x_1^{k_1}\cdots
x_n^{k_n}
\end{gather*}
for any $i=1,\ldots,n-1$, $j=1,\ldots,n$, $\zeta\in \mathbb C^\times$.
Both actions extend the defining action of~$\mathbb G_n$ on~$V$.

\item[$(b)$] The action $\mathop{\triangleright}\nolimits_+$ of $\mathbb G_n$ on $S(V)$ is compatible with the natural
commutative multiplication on~$S(V)$, in the sense that~$\mathbb G_n$ acts by automorphisms of the algebra~$S(V)$.

\item[$(c)$] The action $\mathop{\triangleright}\nolimits_-$ of $\mathbb G_n$ on $S(V)$ is compatible    with the
algebra structure $S_{\mathbf{-1}}(V)$ on $S(V)$ $($which is $S_{\mathbf q}(V)$ with $q_{ij}=-1$ for all $i\ne j)$;

\item[$(d)$] $\rho_+(\mathbb C \mathbb G_n)= \rho_-(\mathbb C \mathbb G_n)$, where $\rho_\pm\colon \mathbb C\mathbb G_n \to
\End_\mathbb C S(V)$ are the algebra homomorphisms arising from the actions $\mathop{\triangleright}\nolimits_\pm$.
Moreover, $\rho_-=\rho_+ \circ J$ for some algebra automorphism~$J$ of~$\mathbb C \mathbb G_n$.
\end{enumerate}
\end{Proposition}

\begin{Remark}
\label{rem:injectivity}
By a~powerful result on group actions on integral domains, see Corollary~\ref{cor:premet-dedekind} from the
Appendix taken with $R=\mathbb C$, $A=S(V)$, the action $\mathop{\triangleright}\nolimits_+$ is faithful, and
moreover the corresponding algebra homomorphism $\rho_+\colon \mathbb C \mathbb G_n \to \End_\mathbb C S(V)$ is
injective.
It follows from (d) that $\rho_-\colon \mathbb C \mathbb G_n\to \End_\mathbb C S_{\mathbf{-1}}(V)$ is also injective.
This fact does not readily follow from classical results.
\end{Remark}



At this point, we restrict our attention to f\/inite subgroups~$G$ of $\mathbb G_n$~-- namely, to Shephard--Todd's
imprimitive complex ref\/lection groups $G(m,p,n)$ and the groups $W_{{\mathcal C},{\mathcal C}'}$, introduced
independently in~\cite{BBselecta} and~\cite{KKZ} and def\/ined as follows.
Let $n\ge 1$ and ${\mathcal C}'\subseteq{\mathcal C}$ be two f\/inite subgroups of $\mathbb C^\times$ of orders $\frac
mp$, $m$ respectively.
Then
\begin{gather*}
G(m,p,n)=\{tw\in {\mathcal C}^n\rtimes \mathbb S_n: \det t\in{\mathcal C}'\},
\\
W_{{\mathcal C},{\mathcal C}'} = \{ tw\in {\mathcal C}^n\rtimes \mathbb S_n: \det(tw)\in {\mathcal C}'\}.
\end{gather*}
The similarity of the two def\/initions manifests itself in our f\/irst main result  where a~correspondence, $\mu$,
between the two classes of subgroups of $\mathbb G_n$ is established.

\begin{maintheorem}
\label{thm:main}
Given $n\ge 1$, an even $m\ge 2$ and a~divisor $p$ of $m$, let $G=G(m,p,n)$.
Then there exists a~unique finite subgroup $ \mu(G) \subset \mathbb G_n $ such that the restriction of
$\mathop{\triangleright}\nolimits_-$ onto $\mu(G)$ is mystically equivalent to $\mathop{\triangleright}\nolimits_+$ on~$G$.
In fact,
\begin{gather*}
\mu(G)=W_{\mathcal C,\mathcal C'},
\end{gather*}
where $|{\mathcal C}|=m$, $|{\mathcal C}'|=\frac mp$.
\end{maintheorem}

The def\/inition of the group $\mu(G)$ suggests that the invariants of $\mu(G(m,p,n))=W_{{\mathcal C},{\mathcal C}'}$
should be viewed in the noncommutative algebra $S_{\mathbf{-1}}(V)$, where this group acts via
$\mathop{\triangleright}\nolimits_-$.
To describe these invariants, introduce the following elements in the space $S(V)$:
\begin{gather*}
p_k^{(m)} = \sum\limits_{i=1}^n x_i^{km},
\qquad
k=1,\ldots,n-1,
\qquad
r^{(l)} = x_1^{l}x_2^{l}\cdots x_n^{l}.
\end{gather*}
The classical result of Shephard--Todd and Chevalley asserts that the subalgebra $S(V)^{G(m,p,n)}$ of~$S(V)$ (with
respect to the natural commutative product on~$S(V)$) equals
the polynomial algebra $\mathbb C\big[p_1^{(m)},\ldots,p_{n-1}^{(m)}, r^{(\frac{m}{p})}\big]$.

Our mystic equivalence construction immediately leads to the following result, f\/irst obtained by Kirkman, Kuzmanovich
and Zhang as a~key ingredient in the classif\/ication theorem~\cite[Theorem~1.1]{KKZ}.

\begin{Theorem}
In the notation of Theorem~{\rm \ref{thm:main}}, let $m$ be even.
Then in the algebra $S_{\mathbf{-1}}(V)$,
\begin{enumerate}\itemsep=0pt
\item[$(a)$] the elements $p_1^{(|\mathcal C|)},\ldots,p_{n-1}^{(|\mathcal C|)}, r^{(|\mathcal C'|)}$ are pairwise commuting
invariants of the group $W_{\mathcal C,\mathcal C'}$;

\item[$(b)$] $S_{\mathbf{-1}}(V)^{W_{\mathcal C,\mathcal C'}} = \mathbb C\big[p_1^{(|\mathcal C|)},\ldots,p_{n-1}^{(|\mathcal C|)},r^{(|\mathcal C'|)}\big]$.
\end{enumerate}
\end{Theorem}

\begin{Remark}%\label{rem:EKL}
This result shows that the condition that $m$ is even in Theorem~\ref{thm:main} is important: the symmetric group
$\mathbb S_n=G(1,1,n)$ does not have a~mystical counterpart, and indeed the correspondence $\mu$ cannot be extended to
groups $G(m,p,n)$ where $m$ is odd.
This happens because the space of the invariants of $G(m,p,n)$ in $S(V)$ is not closed under the multiplication in
$S_{\mathbf{-1}}(V)$, and therefore cannot be the space of invariants of a~group acting by automorphisms of
$S_{\mathbf{-1}}(V)$.
\end{Remark}

The groups $G=G(m,p,n)$ and $\mu(G)$ are of the same order $\frac{m^n n!}{p}$, and, informally, they ``look very similar''.
Our next main result makes this informal statement more precise.
We keep the notation from Theorem~\ref{thm:main} and denote by $R$ the ring $\mathbb
Z\big[\frac{1+\mathbf{i}}{2}\big]\subset\mathbb C$.
\begin{maintheorem}
\label{thm:main2}
For all~$G$ as in Theorem~{\rm \ref{thm:main}}, the group rings $RG$ and $R\mu(G)$ are isomorphic.
In particular, $ \mathbb C G \cong \mathbb C \mu(G) $.
\end{maintheorem}
\setcounter{section}{2}
\setcounter{subsection}{0}
\setcounter{subsubsection}{0}
\setcounter{Theorem}{14}
\setcounter{equation}{0}
% Original source lines 390--616
\section{Proofs of results from Section~\ref{sect:main}}

We will repeatedly use the following straightforward technical fact about the root system of type~$A_{n-1}$, the proof
of which is left to the reader as an exercise.
\begin{claim}
\label{claim:cocycle}
Let $A$ be a~multiplicatively written abelian group.
Let $\phi_{ij}\colon \mathbb Z^n\to A$, $1\le i,j\le n$, $i\ne j$, be a~system of maps satisfying
\begin{gather*}
\phi_{ij}(w(\mathbf k)) = \phi_{w^{-1}(i),w^{-1}(j)}(\mathbf k),
\qquad
\phi_{ji}(\mathbf k)\phi_{ij}(\mathbf k)=1
\end{gather*}
for all $\mathbf k=(k_1,\ldots,k_n)\in\mathbb Z^n$, $w\in\mathbb S_n$, where $w(\mathbf k)$ stands for
$(k_{w^{-1}(1)},\ldots,k_{w^{-1}(n)})$.
Denote
\begin{gather*}
\phi_w(\mathbf k) = \prod\limits_{i<j,w(i)>w(j)} \phi_{ij}(\mathbf k).
\end{gather*}
Then:
\begin{enumerate}\itemsep=0pt
\item[$(a)$] For all $w',w\in \mathbb S_n$,
\begin{gather*}
\phi_{w'w}(\mathbf k) = \phi_{w'}(w(\mathbf k))\phi_w(\mathbf k).
\end{gather*}

\item[$(b)$] If for all $i\ne j$ and for all $\mathbf k,\mathbf{k'}\in \mathbb Z^n$ one has
\begin{gather*}
\phi_{ij}(\mathbf k+\mathbf{k'}) = a_{ij}^{k_ik_j'- k_i'k_j} \phi_{ij}(\mathbf k)\phi_{ij}(\mathbf{k'}),
\end{gather*}
then for all $w\in \mathbb S_n$ one has
\begin{gather*}
\langle \mathbf k,\mathbf{k'}\rangle \phi_{w}(\mathbf k+\mathbf{k'}) = \langle w(\mathbf k),w(\mathbf{k'})\rangle
\phi_{w}(\mathbf k)\phi_{w}(\mathbf{k'}),
\end{gather*}
where $\langle \mathbf k,\mathbf{k'}\rangle= \prod\limits_{i<j}a_{ij}^{k_i'k_j}$.
Here $a_{ij}$ are elements of $A$ such that $a_{ji}=a_{ij}$.
\end{enumerate}
\end{claim}

\subsection*{Proof of Proposition~\ref{prop:automorphism}}

Observe that the group $\mathbb G_n\subset \mathrm{GL}_n(\mathbb C)$ is generated by $\mathbb S_n$ and $(\mathbb
C^\times)^n$ subject to the semidirect product relations
\begin{gather*}
wt_j^{(\zeta)}=t_{w(j)}^{(\zeta)} w,
\qquad
\text{for all} \ \ w\in \mathbb S_n,
\quad
j\in\{1,\ldots,n\},
\quad
\zeta\in\mathbb C^\times.
\end{gather*}
In what follows, we use the abbreviation $x^{\mathbf k}$ to denote $x_1^{k_1}\cdots x_n^{k_n}$, where $\mathbf
k=(k_1,\ldots,k_n)\in \mathbb Z_{\ge 0}^n$.
For $c\in \mathbb C^\times$, def\/ine $\phi_{ij}^{(c)}\colon \mathbb Z^n\to \mathbb C^\times$ by
\begin{gather*}
\phi_{ij}^{(c)}(\mathbf k) = (-1)^{k_ik_j} c^{\overline k_i-\overline k_j},
\qquad
\text{where}
\qquad
\overline k=(k\ \mathrm{mod}\ 2)\in\{0,1\}.
\end{gather*}
Additionally, def\/ine $\phi_{ij}^{(0)}(\mathbf k)=1$ for all $\mathbf k$.
Because $\phi_{ij}^{(c^{-1})}(\mathbf k)=\phi_{ij}^{(c)}(\mathbf k)^{-1}=\phi_{ji}^{(c)}(\mathbf k)$, the system~$\phi_{ij}^{(c)}$ satisf\/ies the condition in Claim~\ref{claim:cocycle} and hence gives rise to functions
$\phi_w^{(c)}$.
The following lemma is then immediate from Claim~\ref{claim:cocycle}(a).

\begin{Lemma}
\label{lem:action}
For each $c\in\mathbb C$, the formula
\begin{gather*}
w t_1^{(\zeta_1)}\cdots t_n^{(\zeta_n)} \mathop{\triangleright}\nolimits_c x^{\mathbf k} = \phi^{(c)}_w(\mathbf k)
\left(\prod\limits_{j=1}^n \zeta_j^{k_j}\right) x^{w(\mathbf k)}
\end{gather*}
defines an action $\mathop{\triangleright}\nolimits_c$ of the group $\mathbb G_n$ on the space $S(V)$.
\end{Lemma}

The actions $\mathop{\triangleright}\nolimits_0$, $\mathop{\triangleright}\nolimits_1$ given by Lemma~\ref{lem:action}
coincide on the generators $s_i$, $t_j^{(\zeta)}$ of $\mathbb G_n$ with the actions
$\mathop{\triangleright}\nolimits_+$, $\mathop{\triangleright}\nolimits_-$ def\/ined in part~(a) of
Proposition~\ref{prop:automorphism}.
Hence part~(a) of the proposition is proved.



Denote by $\bullet_+$, respectively $\bullet_-$, the multiplication on the algebra $S(V)$, respectively
$S_{\mathbf{-1}}(V)$.
One has the following multiplication rule for monomials:
\begin{gather*}
x^{\mathbf k}\bullet_\pm x^{\mathbf{k'}} = \langle \mathbf k,\mathbf k'\rangle_\pm x^{\mathbf{k}+\mathbf{k'}}
\qquad
\text{for all}\quad\mathbf{k}, \mathbf{k'}\in \mathbb Z_{\ge 0}^n,
\end{gather*}
where $\langle \mathbf k,\mathbf k'\rangle_\pm$ is as given in Claim~\ref{claim:cocycle}(b) with all $a_{ij}=\pm1$, $i\ne j$.

It is enough to check that $w\mathop{\triangleright}\nolimits_\pm$ and
$t_j^{(\zeta)}\mathop{\triangleright}\nolimits_\pm$ are automorphisms of the respective algebra structures on $S(V)$.
We apply these actions to both sides of the multiplication rule for monomials and check that the results are equal.
This is trivial for $t_j^{(\zeta)}\mathop{\triangleright}\nolimits_\pm$.
For $w\mathop{\triangleright}\nolimits_\pm$ where $w\in \mathbb S_n$,   the equality is guaranteed by
Claim~\ref{claim:cocycle}(b) and Lemma~\ref{lem:action}, applied to functions $\phi_{ij}^+=\phi_{ij}^{(0)}$,
respectively $\phi_{ij}^-=\phi_{ij}^{(1)}$.
This proves parts (b), (c) of Proposition~\ref{prop:automorphism}.

Now let us prove part (d) of Proposition~\ref{prop:automorphism}.
Clearly, the natural $\mathbb S_n$-action on the group $(\mathbb C^\times)^ n\subset \mathbb G_n$ extends to that on the
group algebra $\mathbb C (\mathbb C^\times)^n$.

For each $c\in \mathbb C\setminus\{0\}$, $w\in \mathbb S_n$, def\/ine the element $Q_w^{(c)}\in \mathbb C (\mathbb C^\times)^n$ by
\begin{gather*}
Q_w^{(c)} = \prod\limits_{i<j:w(i)>w(j)} Q_{ij}^{(c)},
\end{gather*}
where $ Q_{ij}^{(c)} = \frac14 \big(\big(c+c^{-1}\big)\big(1-\g i\g j\big)+ \big(c-c^{-1}+2\big)\g i+\big(c^{-1}-c+2\big)\g j\big)$, $1\le i,j\le n$.
\begin{Lemma}\label{lem:Q}
For all $w',w\in \mathbb S_n$, $c\in\mathbb C^\times$, $Q_w^{(c)}Q_w^{(c^{-1})}=1$ and
$Q_{w'w}^{(c)}=w^{-1}(Q_{w'}^{(c)})\cdot Q_w^{(c)}$.
\end{Lemma}
\begin{proof}
Denote by $\mathcal F$ the algebra of all functions from $\mathbb Z^n$ to $\mathbb C$ with pointwise addition and
multiplication.
Clearly, the assignment $t_1^{(\zeta_1)}\ldots t_n^{(\zeta_n)} \mapsto (\mathbf k\mapsto \zeta_1^{k_1}\ldots
\zeta_n^{k_n})$ def\/ines a~group homomorphism $(\mathbb C^\times)^n\to \mathcal F^\times $ which extends to an injective
map $\Psi\colon \mathbb C(\mathbb C^\times)^n\hookrightarrow \mathcal F$.

It is easy to check that $\Psi(Q_{ij}^{(c)})=\phi_{ij}^{(c)}$ where $\phi_{ij}^{(c)}$ is as in Lemma~\ref{lem:action}.
In particular, $Q_{ij}^{(c)}Q_{ij}^{(c^{-1})}=1$ since $\phi_{ij}^{(c)}\phi_{ij}^{(c^{-1})}=1$ and $\Psi$ is injective.
This proves the f\/irst assertion of the lemma.
To prove the second assertion, apply $\Psi^{-1}$ to Claim~\ref{claim:cocycle}(a) and use the fact that the function
$\mathbf k\mapsto \phi(w(\mathbf k))$ is mapped by $\Psi^{-1}$ to $w^{-1}(\Psi^{-1}(\phi))$ for all functions $\phi$
from the subgroup of $\mathcal F^\times$ generated by $\{\phi_{ij}^{(c)}:i\ne j\}$.
\end{proof}

Now for each $c\in\mathbb C^\times$ def\/ine the $\mathbb C$-linear map $J_c\colon \mathbb C\mathbb G_n\to \mathbb
C\mathbb G_n$ by the formula
\begin{gather*}
J_c(wt)=wtQ^{(c)}_w,
\qquad
t\in \mathbb C(\mathbb C^\times)^n,
\qquad
w\in \mathbb S_n.
\end{gather*}
\begin{Lemma}\label{lem:J_c}
For each $n\ge 1$,
\begin{enumerate}\itemsep=0pt
\item[$(a)$] $J_c$ is an algebra automorphism of $\mathbb C \mathbb G_n$ with inverse $J_{c^{-1}}$.

\item[$(b)$] $\rho_+\circ J_c=\rho_c$, where $\rho_c:\mathbb C \mathbb G_n\to \End_\mathbb C S(V)$ is the algebra homomorphism
corresponding to $\mathop{\triangleright}\nolimits_c$.
\end{enumerate}
\end{Lemma}

\begin{proof}
On the one hand, $J_c(w't'w t)=J_c(w'w\cdot w^{-1}(t')t)=w'w \cdot w^{-1}(t')tQ_{w'w}^{(c)}$.
On the other hand,
\begin{gather*}
J_c(w't')J_c(w t) = w't'Q_{w'}^{(c)} wt Q^{(c)}_w=w'w \cdot w^{-1}(t')tw^{-1}(Q_{w'}^{(c)}) Q_w^{(c)}=J_c(w't'wt)
\end{gather*}
by the second assertion of Lemma~\ref{lem:Q}.
Hence $J_c$ is a~homomorphism of algebras.
Now{\samepage
\begin{gather*}
J_c(J_{c^{-1}}(wt))=J_c\big(wtQ_w^{(c^{-1})}\big)=wtQ_w^{(c^{-1})}Q_w^{(c)}=wt
\end{gather*}
by the f\/irst assertion of Lemma~\ref{lem:Q}.
This proves part (a) of the lemma.}

Prove (b).
In view of Lemma~\ref{lem:action}, it suf\/f\/ices to show that $Q_w^{(c)}\mathop{\triangleright}\nolimits_+ x^{\mathbf
k}=\phi_w^{(c)}(\mathbf k) x^{\mathbf k}$.
Indeed,
\begin{gather*}
Q_w^{(c)}\mathop{\triangleright}\nolimits_+ x^{\mathbf k}
= \frac14 \big(\big(c+c^{-1}\big)\big(1-(-1)^{k_i+k_j}\big)+\big(c-c^{-1}+2\big)(-1)^{k_i}+\big(c^{-1}-c+2\big)(-1)^{k_j}\big)x^{\mathbf k}
\\
\phantom{Q_w^{(c)}\mathop{\triangleright}\nolimits_+ x^{\mathbf k}}
=\phi_{ij}^{(c)}(\mathbf k)x^{\mathbf k}.
\end{gather*}

Finally,
\begin{gather*}
 J_c(wt)\mathop{\triangleright}\nolimits_+ x^{\mathbf k}
=\big(wtQ_w^{(c)}\big)\mathop{\triangleright}\nolimits_+x^{\mathbf k}
=(wt)\mathop{\triangleright}\nolimits_+\big(Q_w^{(c)}\mathop{\triangleright}\nolimits_+x^{\mathbf k}\big)
=\phi_w^{(c)}(\mathbf k)(wt)\mathop{\triangleright}\nolimits_+x^{\mathbf k} =wt\mathop{\triangleright}\nolimits_c x^{\mathbf k}. \!\!\!\tag*{\qed}
\end{gather*}
\renewcommand{\qed}{}
\end{proof}

Taking $c=1$ in Lemma~\ref{lem:J_c}(b), we settle Proposition~\ref{prop:automorphism}(d).
Proposition~\ref{prop:automorphism} is proved.


\subsection*{Proof of Theorem~\ref{thm:main}} We retain the notation from the proof of
Proposition~\ref{prop:automorphism}.
Let $G=G(m,p,n)$ as in the theorem, and denote $T:=G\cap (\mathbb C^\times)^n$.
Let $e_T=\sum\limits_{t\in T}t\in \mathbb C T\subset \mathbb C G$.
Clearly, $\g ie_T=\g 1e_T$ for all $i=1,\ldots,n$.
Hence $Q_{ij}^{(c)}e_T=\g1e_T$ for all $c\in\mathbb C^\times$, so that $Q_w^{(c)}e_T=t_1^{(\det w)}e_T$ for all $w\in
\mathbb S_n$.
Since, as sets, $G=\{wt \,|\, w\in \mathbb S_n, t\in T\}$, one has $e_G=\sum\limits_{w\in\mathbb S_n}w\cdot e_T$ and
\begin{gather*}
J_c(e_G)= \sum\limits_{w\in \mathbb S_n} J_c(w)e_T = \sum\limits_{w\in \mathbb S_n} wQ_w^{(c)}e_T =\sum\limits_{w\in
\mathbb S_n} wt_1^{(\det w)}e_T = e_{\mu(G)},
\end{gather*}
because, as sets, $\mu(G)=\big\{wt_1^{(\det w)} t \,|\, w\in \mathbb S_n, t\in T\big\}$.
Since a~subgroup $G'$ of $\mathbb G_n$ is uniquely determined by $e_{G'}\in\mathbb C \mathbb G_n$, the group $\mu(G)$ is
a~unique subgroup $G'$ of $\mathbb G_n$ such that $J_c(e_G)=e_{G'}$.

Finally, by Lemma~\ref{lem:J_c}(b), $\rho_c(e_G)=\rho_+(e_{\mu(G)})$.
Setting $c=1$ and using injectivity of $\rho_+$, see Remark~\ref{rem:injectivity}, completes the proof of
Theorem~\ref{thm:main}.


\subsection*{Proof of Theorem~\ref{thm:main2}} Let $G=G(m,p,n)$.
Consider the restriction of $J_\mathbf{i}$, $\mathbf{i}=\sqrt{-1}$ to $R G=\mathbb S_n\cdot R T$ where $T=G\cap (\mathbb
C^\times)^n=\mu(G)\cap (\mathbb C^\times)^n$.
Observe that $J_\mathbf{i}(T)=T$ and
\begin{gather*}
J_\mathbf{i}(s_i)=s_iQ_{s_i}^{(\mathbf{i})}=s_i Q_{i,i+1}^{(\mathbf{i})}
=\frac12 s_i\left((1+\mathbf{i})t_i^{(-1)}+(1-\mathbf{i})t_{i+1}^{(-1)}\right)
=\frac{1+\mathbf{i}}{2}\sigma_i+\frac{1-\mathbf{i}}{2}\sigma_i^{-1},
\end{gather*}
where $\sigma_i=s_i t_i^{(-1)}\in \mu(G)$.
Thus, $J_\mathbf{i}(R G)\subseteq R\mu(G)$, so that the automorphism $J_\mathbf{i}$ of $\mathbb C \mathbb G_n$ restricts
to an isomorphism $R G \xrightarrow{\sim} R \mu(G)$.
Theorem~\ref{thm:main2} is proved.
\setcounter{section}{3}
\setcounter{subsection}{0}
\setcounter{subsubsection}{0}
\setcounter{Theorem}{11}
\setcounter{equation}{0}
% Original source lines 827--918
\section{Appendix}

The aim of this section is to prove the following important result about group ring actions.

\begin{Theorem}
\label{thm:Premet}

Let $A$ be an integral domain and~$G$ be a~group of ring automorphisms of $A$.
Then:
\begin{enumerate}\itemsep=0pt
\item[$(a)$] the natural map $\rho\colon AG \to \End_\mathbb Z(A)$ given by $\big(\rho\big(\sum\limits_i a_i g_i\big)\big)(a)=\sum\limits_i a_i g_i(a)$ is injective;
\item[$(b)$] with respect to the natural ring structure on $\End_\mathbb Z(A)$ and the semidirect product structure $A\rtimes
\mathbb Z G$ on $AG$, the map $\rho$ is a~ring homomorphism.
\end{enumerate}
\end{Theorem}

The following is immediate from the theorem.

\begin{Corollary}
\label{cor:premet-dedekind}
In the notation of the theorem, let $R$ be a~subring of $A^G$.
Then the restriction of $\rho$ to $RG$ is an injective ring homomorphism
\begin{gather*}
RG\hookrightarrow \End_R A,
\end{gather*}
where $RG$ is the ordinary group ring of~$G$.
\end{Corollary}

\begin{proof}[Proof of Theorem~\ref{thm:Premet}]
We need the following generalization of the celebrated Dedekind's lemma.
\begin{Lemma}
\label{lem:dedekind}
Let $A$ be an integral domain and $B$ be a~multiplicative monoid.
Let $\mathcal G$ be a~set of monoid homomorphisms $B\to A$.
Then the natural $A$-linear map
\begin{gather*}
\rho\colon \ A\mathcal G \to {\rm Fun}(B,A),
\qquad
\left(\rho\left(\sum\limits_g a_g g\right)\right)(b) = \sum\limits_g a_g g(b)
\end{gather*}
is injective.
$($Here $A\mathcal G=\oplus_{g\in \mathcal G} Ag$ is the free $A$-module generated by $\mathcal G.)$
\end{Lemma}
\begin{proof}
Assume for contradiction that $\rho$ is not injective.
Let $\mathcal G _0$ be a~minimal f\/inite subset of $\mathcal G $ such that there exists a~non-zero element
$\sum\limits_{g\in \mathcal G _0}a_g g$ in the kernel of $\rho$.
That is,
\begin{gather*}
\sum\limits_{g\in \mathcal G _0}a_g g(b)=0
\qquad
\text{for all} \ \ b\in B.
\end{gather*}
Clearly, all $a_g$ are non-zero and $|\mathcal G _0|>1$ because $A$ has no zero divisors.
In particular, for each $b,b'\in B$,
\begin{gather*}
\sum\limits_{g\in \mathcal G _0} a_g g(b'b)=\sum\limits_{g\in \mathcal G _0} a_g g(b')g(b)=0,
\end{gather*}
because all $g\in \mathcal G _0$ are homomorphisms from $B$ to $A$.
Furthermore, f\/ix $h\in \mathcal G _0$.
Combining the above identities, we obtain
\begin{gather*}
\sum\limits_{g\in \mathcal G _0\setminus\{h\}} a_g g(b') (g(b) - h(b)) = 0
\qquad
\text{for all}\quad
b,b'\in B.
\end{gather*}
That is, for each $b\in B$, the element
\begin{gather*}
k_b:=\sum\limits_{g\in \mathcal G _0\setminus\{h\}} a_g (g(b) - h(b)) g
\end{gather*}
belongs to $\ker \rho$.
The minimality of $\mathcal G _0$ implies that $k_b=0$ for all $b\in B$, which is equivalent to $g(b)=h(b)$ for all
$b\in B$ and all $g\in \mathcal G _0$.
That is, $|\mathcal G _0|=1$, a~contradiction.
\end{proof}

We now use Lemma~\ref{lem:dedekind} for $A$ as in Theorem~\ref{thm:Premet} with $B=A\setminus\{0\}$ and $\mathcal G=G$
viewed as homomorphisms of monoids $A\setminus\{0\} \to A$.
Since $\End_\mathbb Z A$ is naturally a~subset of ${\rm Fun}(A\setminus\{0\},A)$ and $\rho(AG)\in \End_\mathbb Z
A\subset{\rm Fun}(A\setminus\{0\},A)$, part (a) of Theorem~\ref{thm:Premet} is proved.

To prove (b), note that $\End_\mathbb Z A$ is naturally a~ring, with multiplication being composition of maps.
The semidirect product multiplication on $A G$ is given by the formula $(ag)(a'g')=ag(a')\cdot gg'$ for all $a,a'\in A$,
$g,g'\in G$.
Then
\begin{gather*}
\rho((ag)(a'g'))(b)=a\cdot g(a'\cdot g'(b))=a\cdot g(a')\cdot g(g'(b)) = ag(a') (gg')(b)
\end{gather*}
for all $a,a',b\in A$, $g,g'\in G$.
Part (b) of the theorem is proved.
\end{proof}
\begin{thebibliography}{99}
\bibitem[1]{BBselecta}
Bazlov Y., Berenstein A., Noncommutative {D}unkl operators and braided
  {C}herednik algebras, \href{http://dx.doi.org/10.1007/s00029-009-0525-x}{\textit{Selecta Math.~(N.S.)}} \textbf{14} (2009),
  325--372, \href{http://arxiv.org/abs/0806.0867}{arXiv:0806.0867}.

\bibitem[2]{KKZ}
Kirkman E., Kuzmanovich J., Zhang J.J., Shephard--{T}odd--{C}hevalley theorem
  for skew polynomial rings, \href{http://dx.doi.org/10.1007/s10468-008-9109-2}{\textit{Algebr. Represent. Theory}} \textbf{13}
  (2010), 127--158, \href{http://arxiv.org/abs/0806.3210}{arXiv:0806.3210}.

\end{thebibliography}
\end{document}
